\chapter{Conditionnement, indépendance et Bayes}
\section*{Exercice 19.1 -- Trois probabilités}
Les définitions donnent $P_A(B)=0{,}3/0{,}4=0{,}75$ et $P_B(A)=0{,}3/0{,}5=0{,}6$. La réunion vaut $0{,}4+0{,}5-0{,}3=0{,}6$. Le conditionnement change le dénominateur.

\section*{Exercice 19.2 -- Probabilités totales}
Les origines A et B partitionnent la production. Donc $P(D)=0{,}6(0{,}02)+0{,}4(0{,}05)=0{,}032$, soit $3{,}2\%$. La moyenne simple des taux serait incorrecte car les poids des machines diffèrent.

\section*{Exercice 19.3 -- Bayes}
$P(B\cap D)=0{,}4(0{,}05)=0{,}02$. Donc $P_D(B)=0{,}02/0{,}032=5/8=62{,}5\%$. Une machine produisant moins de pièces peut être majoritaire parmi les pièces défectueuses.

\section*{Exercice 19.4 -- Tester l'indépendance}
Le produit $(1/3)(3/5)=1/5$ égale la probabilité de l'intersection : ils sont indépendants. Le théorème sur les complémentaires donne aussi leur indépendance. Directement, $P(\overline A\cap\overline B)=1-(1/3+3/5-1/5)=4/15=(2/3)(2/5)$.

\section*{Exercice 19.5 -- Incompatibilité}
L'intersection vide a une probabilité nulle. Mais $P(A)P(B)>0$. L'égalité définissant l'indépendance est donc impossible. Savoir que $A$ se produit exclut entièrement $B$.

\section*{Exercice 19.6 -- Événement nul}
Le quotient $P(A\cap B)/P(A)$ est indéfini car son dénominateur est nul. En revanche, $A\cap B\subset A$ implique $P(A\cap B)=0=P(A)P(B)$ : l'indépendance, définie par le produit, reste vraie. Ne pas employer une équivalence conditionnelle hors de son domaine.
